% §6.7 Cross-method orthogonality between Shen, Baur, and counterfactual.
% Numbers locked against results/replications/cross_method_concordance.csv
% and hotelling_t2_cross_method.csv (both refreshed with the clean CF run).
% Spearman and sign-agreement computed on the per-city point estimates.
% Clean k=12, no Portland exclusion. The v2 "joint factor model" subsection
% (lambda/sigma values, 12/12 credible intervals) had no backing file and is
% removed; a latent one-factor pooling is noted as planned future work only.

\subsection{Cross-method orthogonality}
\label{sec:cross-method-12}

Sections~\ref{sec:shen-12city} and~\ref{sec:baur-12city} each identify
a different small subset of metropolitan markets in which a
text-on-price effect can be detected at conventional significance. The
natural question is whether the channels agree on which markets carry
an effect, or whether they capture substantively different aspects of
listing language. We answer it by computing per-market rank and sign
correlations between the Shen-Ross uniqueness DML estimate and the
Baur PC1 DML estimate, and we add the counterfactual rewriting total
effect from Section~\ref{sec:counterfactual-12} as a third axis.

\paragraph{Pairwise rank correlations.}
Table~\ref{tab:cross-method-12} reports the per-market point estimates
for the three channels. Across the twelve markets the Spearman rank
correlation between the Shen-Ross estimate and the counterfactual
total effect is $+0.64$ ($p = 0.02$), the strongest of the three
pairings, while the correlation between the Shen-Ross estimate and the
Baur estimate is $-0.59$ ($p = 0.05$) and the correlation between the
Baur estimate and the counterfactual total effect is $-0.52$
($p = 0.09$). With only twelve markets the rank statistic carries a
standard error near $1/\sqrt{12} \approx 0.29$, so we read these
coefficients as descriptive of the cross-market pattern rather than as
precise estimates of agreement. The pattern that emerges is that the
Shen uniqueness channel and the counterfactual-rewriting channel rank
together across markets, whereas the Baur PC1 axis ranks oppositely to
both, which is the direction expected from its negative pooled loading
on the utilitarian-vocabulary axis documented in
Section~\ref{sec:baur-12city}: a market with a strongly positive
Shen uniqueness premium tends to sit low on the Baur axis. We do not
read the negative Baur correlations as evidence that the channels
contradict one another, because the channels are measured on axes with
opposite substantive orientation.

\paragraph{Sign agreement.} Counting markets in which two channels
share the sign of their point estimate, the Shen and counterfactual
channels agree in seven of the twelve markets, the strongest pairwise
sign concordance and consistent with their positive rank correlation,
while Shen and Baur agree in three markets and Baur and the
counterfactual agree in two. No market has all three channels sharing
a common sign, which again reflects the opposite orientation of the
Baur axis rather than methodological inconsistency.

\paragraph{Identification overlap.} A more decision-relevant summary
than rank agreement is which markets each channel identifies at the
$95\%$ level, where the per-method bootstrap confidence interval
excludes zero. The Shen channel identifies ten of the twelve markets,
the counterfactual channel nine, and the Baur axis four, namely New
York, Philadelphia, Denver, and Atlanta, all of which also clear the
Shen threshold. Three markets are identified by all three channels
jointly, New York, Philadelphia, and Denver, and eight of the twelve
are identified by at least two channels, with no market left
unidentified by every channel. Among the markets identified by
exactly two channels, four pair the Shen and counterfactual axes, San
Francisco, Washington DC, Phoenix, and Dallas, and one pairs Shen and
Baur, Atlanta, which mirrors the positive Shen-counterfactual rank
correlation at the level of the identified set.

\paragraph{Substantive interpretation under a triangulation design.}
The three identification strategies are not redundant point estimators
of a single scalar effect; they are deliberately chosen to have
orthogonal sources of bias. The Shen-Ross uniqueness measure leverages
within-market lexical rarity and inherits any bias from the
neighborhood-peer-set definition, the Baur sentence-BERT axis
leverages cross-listing semantic geometry and inherits bias from the
embedding model's pretraining distribution, and the LLM counterfactual
rewriting intervention leverages within-listing minimal-edit
substitution and inherits bias from the rewriter's distributional
shift. Following the causal-triangulation framework of
\citet{lawlor2016triangulation} and \citet{hammerton2021causal},
methods whose biases are largely unrelated provide jointly stronger
evidence on a common causal target than methods that share a bias
structure, regardless of whether their per-market point estimates
agree numerically. This is a per-market restatement of the
\citet{heckman2008econometric} position that different estimators
identify different parameters: each method rejects a method-specific
null about a method-specific functional of the language-price
relationship, while the joint test below rejects the global null that
no functional has any effect in a given market. Under this design the
mixed and imprecise pairwise correlations are the discriminant-validity
signal the design anticipates, and the joint rejection pattern is the
convergent signal, which is the canonical signature predicted when
several low-power estimators each identify a different feature of a
partially-identified estimand \citep{manski2003partial, tamer2010partial}:
the joint test inherits power from the union of identifying
restrictions even when each marginal test does not.

\paragraph{Joint-null Hotelling test.} The most powerful summary of
the cross-method evidence asks, for each market, whether the joint null
$H_0: (\theta_{\mathrm{Shen}}, \theta_{\mathrm{Baur}},
\theta_{\mathrm{CF}}) = (0, 0, 0)$ can be rejected against the
two-sided alternative. Treating the three per-market estimates as
asymptotically independent, since the three methods use disjoint
sources of variation in a Doc2Vec uniqueness score, a sentence-BERT
principal component, and an LLM-rewriting differential, we compute the
Hotelling-style statistic $T^2_m = z^2_{\mathrm{Shen},m} +
z^2_{\mathrm{Baur},m} + z^2_{\mathrm{CF},m}$ per market and refer it to
a $\chi^2_3$ asymptotic distribution. The joint null is rejected in all
twelve of the twelve markets at conventional five-percent significance,
and the rejection survives Benjamini-Hochberg false-discovery
correction across the twelve markets, where the largest adjusted
$q$-value is $0.002$. The weakest joint rejection is Boston, with
$T^2 = 14.5$ and $q = 0.002$, and even there the joint null is
comfortably rejected, so the post-2020 panel contains no market that
reads as a clean text-on-price null. Atlanta, which an earlier draft
treated as a clean null, now rejects strongly with $T^2 = 128$ on the
strength of its Shen and Baur $z$-statistics even though its
counterfactual interval covers zero. The rejection of the joint null
in every market despite the mixed and imprecise pairwise rank
correlations is consistent with each method contributing independent
low-power evidence of a non-zero text-on-price signal, even when no
single method attains significance on that market alone.

We favor the reading that text-on-price effects are real but
heterogeneous across metropolitan markets, with the channel that
detects an effect varying by market, while leaving open the question of
which method is preferred when an analyst can deploy only one. A
constructive complement to the joint Hotelling test would treat the
three per-market estimates as noisy indicators of a single latent
market-level text-price effect and pool them through an anchored
one-factor measurement model in the tradition of latent-variable
meta-analysis, recovering a precision-weighted posterior for each
market's common factor. We have not fit that model on the present
panel and report it here only as a planned extension, so the joint
Hotelling test and the identification-overlap counts above remain the
cross-method evidence on which the conclusions of this section rest.
